\section{Representation level under coordinate deletion}\label{tra:projection-deletion}

We turn from the choice of entropy law to the sparse representation input. The first question is how much diagram level can disappear when a coordinate is deleted.

A split into two coordinate halves changes the representation in which a vector is decomposed. The useful quantity here is how much of its original level disappears in that decomposition. For a partition $\lambda\vdash m$, its level is $m-\lambda_1$. We control the loss of level by averaging tuple coordinates while retaining their side assignments. A content formula gives a lower bound for the resulting projection norms; random orientations of the matching pairs give an upper bound. Comparing them supplies a tail estimate for the lost level.

The dense part retains different information: an entropy inequality valid for every joint law supported on compatible row and column moves. Correlations within the column family remain in the entropy budget until the final step. These two estimates give a complete moment recursion below.

Throughout this section $n=2^d$, $d\ge1$, $G_n=S_n$, and $T_n$ is the sweep from Section~\ref{sec:prelim}. Its regular trace is unnormalized, and one sweep corresponds to $d$ physical shuffles.

\subsection{The content formula and coordinate deletion}
\label{tra:projection-deletion:content-and-injection-preliminaries}
We use standard facts from complex representation theory of symmetric groups. To specify these: irreducibles are indexed by Young diagrams \(\lambda\), have dimension \(d_\lambda\) the number of standard tableaux, are self-dual, and each has multiplicity \(d_\lambda\) in the regular representation. The level for \(G_m\) is \(m-\lambda_1\). We will use the Pieri rule (special case of Young diagram induction/restriction, or the Littlewood-Richardson rule): the possible types \(\rho\) of \(G_l\) paired with the trivial representation of \(G_{m-l}\) when restricting a type \(\lambda\) to \(G_l\times G_{m-l}\) are obtained by removing a horizontal strip of size \(m-l\) from \(\lambda\), with multiplicity one. We also use the content formula (equivalently the Young--Jucys--Murphy formula) for the sum of transpositions \cite[Eq.~(2.1), Proposition~5.3, Theorem~5.8]{VershikOkounkov2005}. Thus the element \(\sum_{1\le h<i\le m}(1-(hi))\) acts on type \(\lambda\) by
\begin{equation}
D_m(\lambda)=\binom m2-\sum_{(u,v)\in\lambda}(v-u)
\label{tra:projection-deletion:eq2}
\end{equation}
where \(u,v\) are row and column.

The representation on ordered injections (ordered distinct sites) of length \(l\) contains only types of level at most \(l\); a type of level \(l\) does occur, by Pieri taking \(\rho=(\lambda_2,\lambda_3,\ldots)\). These statements use the realization as functions of a permutation invariant under the right group on \(m-l\) indices, with \(G_m\) acting on the left; the right \(G_l\) acting on the length-\(l\) ordering thus has types given by Pieri in each left isotypic component. In particular, in the realization of level \(l\) by functions of \(l\) distinct sites, summing out any one of the \(l\) sites (others fixed) gives zero since that projection is equivariant with image in functions of \(l-1\) sites.
\begin{proposition}\label{tra:projection-deletion:level}
\label{tra:projection-deletion:level-statement}
There are absolute \(\alpha>0\), \(M_0<\infty\), and \(0<h<1\), with
\begin{equation}
\|T_n\|_\lambda^2\ \le\ h\left(\frac{M_0 k}{n}\right)^{\alpha k}, \qquad k=n-\lambda_1\ge 1,
\label{tra:projection-deletion:eq3}
\end{equation}
where the norm is operator norm on the indicated representation.
\end{proposition}
\begin{proof}\label{tra:projection-deletion:level-proof}
\label{tra:projection-deletion:random-split-deletion-proof}
It suffices to obtain a nontrivial estimate when \(k/n\) is sufficiently small; outside that range, choosing \(M_0\) large will make the right side at least one. For \(k=1\) the norm is zero: on single-site functions the switches average each bit, giving global averaging as the product.

We induct on the number of coordinates. The main step is a tail bound for the level lost after the first matching projection. Once this bound is established, we will check the constants in the norm recursion.

For induction, write \(n=2m\), use the pairs of a dimension for the rightmost projection in \(T_n\), and call their flip subgroup \(K\), with projection \(P_K\). There are two sides of size \(m\), with the sites of each pair in opposite sides. The other factors act separately in these two sides, so that
\[
T_n=(T_m\otimes T_m)P_K,
\]
with the tensor product here in the subgroup algebra of \(G_m\times G_m\). We use level \(k\) realized on the injection tuples \(x=(x_1,\ldots,x_k)\). It suffices to bound the squared norm after the two side operators starting with a \(K\)-invariant unit vector \(f\) of our type there. Norms of tuple functions can use uniform counting probability. Such \(f\) has zero sum in each of its coordinates as explained above.

Restrict \(f\) to the different assignments of the \(k\) indices to sides, with \(j\) indices in the first side for a given assignment. Decompose in types within the two sides, of levels \(s_1,s_2\), where \(s_1\le j\), \(s_2\le k-j\). We get a probability law on these data by squared norms of the orthogonal parts. Under it put \(R=k-s_1-s_2\). We give two controls on this law:
\begin{itemize}
\item \(j\) has subgaussian tails around \(k/2\), in particular \(\Pr(|j-k/2|\ge t)\le 2\exp(-2t^2/k)\), and \(\Pr(j=0\ \text{or}\ k)\le 2^{1-k}\).
\item For \(k/n\) sufficiently small, there is an absolute \(C_0\) such that
\end{itemize}
\begin{equation}
\Pr(R\ge r)\ \le\ (C_0 k/n)^{r/2},\qquad 1\le r\le k .
\label{tra:projection-deletion:eq4}
\end{equation}

The first property follows from \(K\)-invariance and averaging by randomizing the pair orientations for the side assignment. Each tuple then uses some double-occupied pairs contributing one index each in the first side, and singles with independent fair side assignment.

To prove \eqref{tra:projection-deletion:eq4}, consider projections indexed by sets \(I\) of \(r\) indices to be ignored, that is, average the positions of those indices knowing the positions of the other indices and \textbf{the sides of all indices}. Write these as \(L_I\). All these averages use the uniform injection law conditioned on the
complete side-assignment vector. When an index is deleted, its side
remains part of the conditioning and only its position is integrated
out. Successive deletions are therefore conditional expectations onto
nested sigma-algebras: their composition is the single average over the
deleted set. We apply the one-slot bound on each resulting shorter
tuple space, retaining the same two left representation types.

We first record why
\begin{equation}
\sum_{|I|=r}\|L_I f\|^2\ \ge\ 2^{-r}\Pr(R\ge r)
\label{tra:projection-deletion:eq5}
\end{equation}
when \(k/n\) is sufficiently small. Work conditionally on assigned sides. Within one side with a tuple of \(l\le k\) positions and left type \(\mu\) of level \(s\), the sum of the projections which average out one coordinate has eigenvalues
\[
l-\frac{D_m(\mu)-D_l(\rho)}{m-l+1}
\]
on the corresponding right types \(\rho\) described by Pieri. Indeed the single-coordinate average for coordinate \(i\), on functions of a permutation, acts on the right by \((1+\sum_{l<q\le m}(iq))/(m-l+1)\). Sum over \(i\le l\), and use \eqref{tra:projection-deletion:eq2} and invariance on the right under \(G_{m-l}\). Now \(\rho\) contains the diagram \(\bar\mu=(\mu_2,\mu_3,\ldots)\) since \(\mu/\rho\) is a horizontal strip. Thus \(D_l(\rho)\ge D_s(\bar\mu)\) by adding boxes (each addition to size \(v\) to go to size \(v+1\) has content at most \(v\)). Also \(D_m(\mu)-D_s(\bar\mu)=(m-s+1)s\). This gives the projection sum lower bound
\[
(l-s)\left(1-\frac{s}{m-l+1}\right)\ \ge\ (l-s)/2.
\]
It holds also with zero indices and the trivial type. For both sides together, summing over the coordinate to remove thus gives squared norms summing to at least half the excess of current tuple length over total level, times the squared norm before removal. After an averaging we can work on the shorter tuples with their marginal law; the averaging preserves the left types when nonzero. On a component of total level $k-R$, iterating this bound through $r$ successive deletions gives $2^{-r}(R)_r$ times the original squared norm, where $(R)_r=R(R-1)\cdots(R-r+1)$. At each stage the conditional expectation acts on the shorter tuple space, with all side assignments retained. The tower property identifies the result with averaging all removed coordinates, and each set $I$ is counted $r!$ times through its possible deletion orders. After division by $r!$, the lower bound is $2^{-r}\binom Rr$, which is at least $2^{-r}$ for $R\ge r$. Summing the components proves \eqref{tra:projection-deletion:eq5}. All projections in \eqref{tra:projection-deletion:eq5} respect assignments and the left types of the side groups.

We have expressed a large loss of level as a large sum of coordinate-projection norms. We now bound the same sum from above using the matching-pair symmetries of $f$.

For the upper bound on the left hand side of \eqref{tra:projection-deletion:eq5}, we can average it with the side partition randomized by \(K\), because \(f\) is invariant under \(K\). Write \(\epsilon(x_i)\) for the side sign of the position \(x_i\) under the random orientations. Given the retained positions \(z\), the numerator sum in an average over the other positions constrained by a specified vector of side signs \((\eta_i)_{i\in I}\) equals
\begin{equation}
\sum_{x|z} f(x)\prod_{i\in I}\frac{1+\eta_i\epsilon(x_i)}{2}
=2^{-r}\left(\prod_{i\in I}\eta_i\right)\sum_{x|z} f(x)\prod_{i\in I}\epsilon(x_i),
\label{tra:projection-deletion:eq6}
\end{equation}
where \(x|z\) indicates all distinct-site extensions. The equality uses the zero sums. In computing the squared norm, we sum squared absolute numerator sums divided by the number of extensions with the given side constraints, over \(z\) and \((\eta_i)\), with the overall tuple space normalization. These denominators are at least \((m-k)^r\). Thus we can use this lower bound before expanding the averaged squares on the right in \eqref{tra:projection-deletion:eq6}.

Two tuples in that expansion, \(x,y\) agreeing off \(I\), must among their \(I\)-entries together hit each pair an even number of times for the sign expectation not to vanish. We bound the resulting nonnegative majorant on absolute values of \(f\) by a row sum (the matrix is symmetric). For fixed \(x\), classify \(I\) as using \(u\) single entries in pairs and \(v\) doubles, \(u+2v=r\), counting just \(x_I\). The number of choices of \(I\) with these data is at most \(k^{u+v}/(u!\,v!)\): choose \(v\) full pairs in \(x\), then the other entries. For each, the possibilities for \(y_I\) are bounded by \(2^u r! n^v/v!\), since they have to use the same \(u\) singleton pairs, and \(v\) doubles. The sum over sign specifications costs \(2^r\), with a factor \(2^{-2r}\) in the square in \eqref{tra:projection-deletion:eq6}. It follows, using \(r!/(u!v!v!)\le 3^r\), that the left hand side of \eqref{tra:projection-deletion:eq5} is at most
\[
\frac{2^{-r}}{(m-k)^r}
\sum_{u+2v=r}\frac{k^{u+v}}{u!\,v!}\cdot\frac{2^u r! n^v}{v!}\ \le\ C_1^r(k/n)^{r/2}
\]
for an absolute \(C_1\) and \(k/n\) sufficiently small. This proves \eqref{tra:projection-deletion:eq4}.

The lost-level tail \eqref{tra:projection-deletion:eq4} is now proved. It remains to show that its cost is smaller than the contraction supplied by the two child sweeps.

To finish the induction for \eqref{tra:projection-deletion:eq3}, use the inductive bound in each side (factor 1 for level 0). For the squared output norm, divided by the desired bound, this gives the upper bound by the expectation of
\begin{equation}
h^{\#\{i:s_i>0\}-1}
\left(\frac{n}{M_0 k}\right)^{\alpha R}
\exp\left(\alpha\sum_{i=1}^2 s_i\log(2s_i/k)\right),
\label{tra:projection-deletion:eq7}
\end{equation}
where zero terms in the sum are set to zero. The sum in this exponent before multiplying by \(\alpha\) is bounded above by
\begin{equation}
C_2\big((j-k/2)^2/k+R\big)
\label{tra:projection-deletion:eq8}
\end{equation}
for some absolute \(C_2\). For \(s_1=j,s_2=k-j\) this follows by the binary entropy formula or Taylor expansion near \(k/2\), and boundedness away from the middle. If \(j/k\) is in \([1/4,3/4]\) and \(R< k/8\), changing to \(s_i\) changes the terms by at most \(O(R)\); outside these conditions the crude upper bound \(O(k)\) suffices.

Here are details that constants can be fixed in the induction. Take \(h=1/20\), and \(\alpha>0\) sufficiently small, with \(\alpha<1/4\). On \(R=0\) and \(k\ge2\), cases with \(j=0\) or \(k\) contribute to the expectation at most \(2^{1-k}2^{\alpha k}\le 0.6\) by making \(\alpha\) small. The other cases with \(R=0\) contribute at most \(2h\), by \eqref{tra:projection-deletion:eq8} and the subgaussian estimate, which by taking \(\alpha\) small gives \(\mathbb E \exp(2\alpha C_2(j-k/2)^2/k)\le 2\). For \(R=r\ge1\), Cauchy-Schwarz with \eqref{tra:projection-deletion:eq8}, the same subgaussian estimate, and \eqref{tra:projection-deletion:eq4} bound the total contribution summed over \(r\ge1\), for \(M_0\ge1\), by
\[
\frac{\sqrt 2}{h}\sum_{r\ge1}
\left(e^{C_2\alpha}(n/k)^\alpha (C_0 k/n)^{1/4}\right)^r.
\]
Choose \(M_0\) so large that whenever \(h(M_0 k/n)^{\alpha k}<1\), \(k/n\) is small enough to apply all the estimates and make this last contribution at most \(1/5\). Thus \eqref{tra:projection-deletion:eq7} gives the bound required. When \(h(M_0 k/n)^{\alpha k}\ge1\) we need no induction by contractivity, and level 1 was settled already. This proves \eqref{tra:projection-deletion:eq3} (with the recursion starting trivially at a single site).
\end{proof}
\subsection{An entropy bound for arbitrary compatible laws}
\label{tra:projection-deletion:grid-law-definition}
The norm estimate controls low representation levels. For the remaining levels we bound the regular trace using an entropy estimate for an array with \(b\) rows and \(a\) columns (\(ab=n\), \(a\le b\le 2a\)). We will use it for large \(a\). Row switches now mean general row permutations \(A=(A_i)_{1\le i\le b}\), each \(A_i\) a permutation on \(a\); and write \(B=(B_c)_{1\le c\le a}\) for column permutations, each on \(b\). We compose right-to-left, so \(BA\) sends \((i,j)\) to row \(B_{A_i(j)}(i)\) and column \(A_i(j)\).

Call \((A,B)\) valid if, for each \(j\), the values
\[
B_{A_i(j)}(i),\quad 1\le i\le b
\]
are distinct. Consider any probability law \(\nu\) supported on valid configurations. Write divergences \(D\) with uniform reference as follows:
\begin{itemize}
\item \(D_{\rm joint}\) for both families together relative to uniform independent permutations,
\item \(D_{A_i},D_{B_c}\) for individual marginal permutations,
\item \(D_B\) for the joint marginal on all column permutations relative to their uniform product law.
\end{itemize}
\begin{proposition}\label{tra:projection-deletion:entropy}
\label{tra:projection-deletion:entropy-statement}
For every probability law $\nu$ supported on valid configurations, with absolute constants and all sufficiently large \(a\),
\begin{equation}
D_{\rm joint}
\ \ge\ n-O(n a^{-1/16})
 +\left(1-\frac{C_3}{\log a}\right)\left(\sum_i D_{A_i}+\sum_c D_{B_c}\right).
\label{tra:projection-deletion:eq9}
\end{equation}
\end{proposition}
\begin{proof}\label{tra:projection-deletion:entropy-proof}
\label{tra:projection-deletion:light-atoms-and-exceptional-counts}
For clarity, by entropy chain rules,
\begin{equation}
D_{\rm joint}=D_B+\sum_i D_{A_i}+I,\qquad
I=\sum_i H(A_i)-H(A\mid B)
\label{tra:projection-deletion:eq10}
\end{equation}
with \(H\) Shannon entropy. We expose all \(B_c\), then the rows \(A_i\) in random order given by increasing independent uniform priorities \(t_i\in[0,1]\) independent of the configuration. Within each row expose the entries in column index order \(j=1,\ldots,a\). The contribution for a cell to \(I\) is the expected relative entropy of its full conditional law given the information available before exposure (including the order), relative to \(p\), its marginal conditional law based on just the history in the same row. Here both laws are for \(A_i(j)\), using \(\nu\) to form the laws. These expected contributions sum to \(I\).

For this cell let the light set consist of \(c\) with \(p(c)\le a^{-1/2}\); denote it by \(L\), and use \(p_L\) for the conditional probability on \(L\) when it has mass. An allowed set is given by the candidates \(c\) for which \(B_c(i)\) does not appear among the previously exposed rows' outputs \(B_{A_{i'}(j)}(i')\) for this column index \(j\). The expected relative entropy for the cell is at least
\begin{equation}
\mathbb E\left[\mathbf 1_{\{A_i(j)\in L\}}\big(-\log p_L(\text{allowed set})\big)\right].
\label{tra:projection-deletion:eq11}
\end{equation}
Indeed one decomposes relative entropy according to whether the entry is light, then keeps only the divergence (weighted by conditional probability) when light, whose support condition gives the log bound. The log expression for the weight only needs to be evaluated when the entry is light and is then finite almost surely.

Say a realized cell is good if \(A_i(j)\in L\), \(p(L)\ge a^{-1/8}\), and
\begin{equation}
\#\{c:B_c(i)=B_{A_i(j)}(i)\}\le a^{1/4}.
\label{tra:projection-deletion:eq12}
\end{equation}
Condition on the realized configuration for a good cell and on \(t_i\), but not on the priorities of other rows. Each candidate \(c\) giving an output distinct from the actual output is forbidden with probability \(t_i\), since precisely one other row gives that candidate output value in column index \(j\). The mass under \(p_L\) of candidates giving the actual output is at most \(a^{-1/8}\), by \eqref{tra:projection-deletion:eq12} and the bounds defining light and good. Consequently the expected allowed mass is at most \(1-t_i+t_i a^{-1/8}\). Notice that goodness and \(p_L\) depend on data of the configuration but not on the priorities. By nonnegativity in \eqref{tra:projection-deletion:eq11} we can restrict there to good cells, and Jensen's inequality and integration in \(t_i\) now yield the lower bound
\begin{equation}
\Pr(\text{cell good})\left(1-O(a^{-1/16})\right)
\label{tra:projection-deletion:eq13}
\end{equation}
for \eqref{tra:projection-deletion:eq11}. This uses \(\int_0^1 -\log(1-t+t a^{-1/8})\,dt \ge 1-O(a^{-1/16})\).

Each good cell has supplied almost one nat of entropy. To obtain a total saving of almost $n$, we must pay for the cells excluded by the three goodness conditions.

We detail how the expected number of bad cells is bounded by
\begin{equation}
O(n a^{-1/16})+\frac{C_4}{\log a}\left(\sum_i D_{A_i}+D_B\right).
\label{tra:projection-deletion:eq14}
\end{equation}
For heavy actual entries, consider in each row the exposure under its marginal law, relative to drawing a uniform permutation sequentially. Ignore its last at most \(O(a^{3/4})\) cells so that all considered draws have at least \(a^{3/4}\) positions left. Under the conditional uniform reference, the set \(\{c:p(c)>a^{-1/2}\}\) has probability at most \(a^{-1/4}\). For any laws \(q,q_0\) the variational inequality for divergence gives, for an event \(E\) and \(v\ge0\),
\begin{equation}
D(q\|q_0)\ge v q(E)-\log(1+q_0(E)(e^v-1)).
\label{tra:projection-deletion:eq15}
\end{equation}
Applying this conditionally with \(v=(\log a)/8\) and summing by the chain rule gives the desired heavy-entry part of \eqref{tra:projection-deletion:eq14}. Cases with \(p(L)<a^{-1/8}\) and a light actual entry have expectation count at most \(n a^{-1/8}\) without any entropy cost.

It remains to address \eqref{tra:projection-deletion:eq12}. Its exceptional count is the same computed over \(i,c\) instead of \(i,j\), looking at the multiplicity of the value \(B_c(i)\) in row \(i\) of the column-permutation data. For groups of cells with equal such value and multiplicity \(>a^{1/4}\), a constant fraction of each group have already had that value appear at least \(a^{1/4}/2\) times in earlier columns (up to harmless rounding, or use threshold \(a^{1/4}/3\)). Reveal the data of \(B\) row by row in deterministic order, with columns in order within each. Ignore rows near the end with fewer than \(b a^{-1/8}\) entries remaining available in a uniform column permutation. The conditional uniform reference probability that the revealed value has already appeared at least \(a^{1/4}/3\) times within the row is at most \(3a^{3/4}/(b a^{-1/8})\le 3a^{-1/8}\). Apply \eqref{tra:projection-deletion:eq15} with \(v=(\log a)/16\) and the chain rule for \(D_B\). This proves the remaining exceptional count bound, hence \eqref{tra:projection-deletion:eq14}.

Summing \eqref{tra:projection-deletion:eq13} and using \eqref{tra:projection-deletion:eq10}, \eqref{tra:projection-deletion:eq14} proves \eqref{tra:projection-deletion:eq9}, using \(D_B\ge\sum_c D_{B_c}\) once the coefficient on \(D_B\) is positive.
\end{proof}
\begin{corollary}\label{tra:projection-deletion:weighted}
\label{tra:projection-deletion:weighted-statement}
Here is a useful weighted form. Put \(\theta=C_3/\log a\) (increasing the absolute constant if needed), so \(0<\theta<1\) for sufficiently large \(a\). Under uniform independent permutations, for nonnegative functions \(z_i\) of a row permutation and \(w_c\) of a column permutation,
\begin{equation}
\begin{split}
\log\mathbb E_{\rm unif}\left[\mathbf 1_{\rm valid}
       \prod_i z_i(A_i)^2\prod_c w_c(B_c)^2\right]
\ \le\ &-n+O(n a^{-1/16})\\
 &+2\sum_i\log\|z_i\|_{2/(1-\theta)}
    +2\sum_c\log\|w_c\|_{2/(1-\theta)}
\end{split}
\label{tra:projection-deletion:eq16}
\end{equation}
where the norms are under the marginal uniform probabilities.
\end{corollary}
\begin{proof}\label{tra:projection-deletion:weighted-proof}
\label{tra:projection-deletion:weighted-variational-proof}
This is the finite-space entropy--Brascamp--Lieb duality of Carlen and Cordero-Erausquin~\cite[Theorem~2.1]{CarlenCorderoErausquin2009}, applied to the compatibility event. We give the specialization to fix the coefficient of every marginal deficit.

The log integral on the left is given by the entropy variational formula, maximizing the expected log product minus \(D_{\rm joint}\) over allowed laws supported on valid pairs. After \eqref{tra:projection-deletion:eq9} the remaining marginal terms are each bounded by the separate variational formulas with coefficient \(1-\theta\). Zero weights can also be treated by a limit.
\end{proof}
\subsection{The dense trace recursion}
\label{tra:projection-deletion:operator-split}
We now apply the weighted entropy estimate to coefficient densities of positive powers of the two smaller sweeps. The fourth trace turns their product into the validity constraint just analyzed. Split \(d\) into parts as equally as possible for \(n=2^d\), thus taking \(a=2^{\lfloor d/2\rfloor}\), \(b=2^{\lceil d/2\rceil}\). Make the array of cube positions accordingly. In group algebra, \(T_n=T^{\rm row}T^{\rm col}\) where the row factor consists of independent runs with operator \(T_a\) in every row, and the column factor similarly uses \(T_b\). Indeed we just divide the dimensions into two consecutive groups; the grouping and naming is in product-of-operators order here. Put
\[
X=(T^{\rm row})^*T^{\rm row},\qquad Y=T^{\rm col}(T^{\rm col})^*
\]
positive operators. Within the subgroup algebras these are products on separate permutations, with individual operators \(X_i\) and \(Y_c\) of the forms \(T_a^*T_a\) and \(T_b T_b^*\), respectively. For traces or norms on subgroup factors we can use their own regular representations. Positive powers stay in the algebra and are likewise products there.
\begin{proposition}\label{tra:projection-deletion:recursion}
\label{tra:projection-deletion:trace-recursion-statement}
Suppose for some \(p\ge2\) we have
\begin{equation}
\operatorname{Tr}_{G_a}|T_a|^{2p}\le 2,\qquad
\operatorname{Tr}_{G_b}|T_b|^{2p}\le 2.
\label{tra:projection-deletion:eq17}
\end{equation}
Let \(p'=p(1+\theta)\) with \(\theta\) as in \eqref{tra:projection-deletion:eq16}. We claim
\begin{equation}
\log\operatorname{Tr}_{G_n}|T_n|^{2p'}\ \le\ O(n a^{-1/16}+\sqrt n\log n).
\label{tra:projection-deletion:eq18}
\end{equation}
\end{proposition}
\begin{proof}\label{tra:projection-deletion:recursion-proof}
\label{tra:projection-deletion:interpolation-and-fourth-trace}
First, a positive-matrix trace comparison gives
\begin{equation}
\operatorname{Tr}|T_n|^{2p'}\ \le\ 
\operatorname{Tr}\left(X^{p'/2}Y^{p'/2}X^{p'/2}Y^{p'/2}\right).
\label{tra:projection-deletion:eq19}
\end{equation}
Apply Lemma~\ref{lem:positive-power-comparison} with
$A_2=T^{\rm row}$, $A_1=T^{\rm col}$, $P=2p'$ and $q=4$.
Then $X=|A_2|^2$, $Y=|A_1^*|^2$, giving exactly
\eqref{tra:projection-deletion:eq19}.

Write \(x(A)=\prod_i x_i(A_i)\) for the coefficients as a density relative to uniform row permutations of \(X^{p'/2}\); similarly \(y(B)=\prod_c y_c(B_c)\) for \(Y^{p'/2}\). The individual positive line operators are
$X_i^{p'/2}$ and $Y_c^{p'/2}$. Applying
Lemma~\ref{lem:coefficient-compatibility} to these operators, before
normalizing their squared coefficient masses, gives
\begin{equation}
\operatorname{Tr}|T_n|^{2p'}\ \le\
\frac{n!}{(a!)^b (b!)^a}
  \ \mathbb E_{\rm unif}\left[\mathbf 1_{\rm valid}|x(A)|^2 |y(B)|^2\right].
\label{tra:projection-deletion:eq20}
\end{equation}

For the line norms we use Lemma~\ref{lem:inverse-fourier-bound}
in each subgroup's own regular representation.
With \(q=2/(1+\theta)\), the Schatten \(q\) norm of \(X_i^{p'/2}\) in its own regular representation is at most \(2^{1/q}\) by \eqref{tra:projection-deletion:eq17}, and similarly for \(Y_c^{p'/2}\). Consequently the individual densities in \eqref{tra:projection-deletion:eq20} have the needed \(2/(1-\theta)\)-norms bounded by 2. Apply \eqref{tra:projection-deletion:eq16} to their absolute values. Since by Stirling's estimate
\[
\log\left(\frac{n!}{(a!)^b(b!)^a}\right)=n+O(\sqrt n\log n),
\]
\eqref{tra:projection-deletion:eq20} proves \eqref{tra:projection-deletion:eq18}.
\end{proof}
\begin{theorem}\label{tra:projection-deletion:moment}
There are real numbers $p_n\ge2$, indexed by dyadic $n\ge2$, with $\sup_n p_n<\infty$ such that
\label{tra:projection-deletion:main-statement}
\begin{equation}
\operatorname{Tr}_{G_n}|T_n|^{2p_n}\ \le\ 1+\tfrac{1}{8}.
\label{tra:projection-deletion:eq1}
\end{equation}
\end{theorem}
\begin{proof}[Proof of Theorem~\ref{tra:projection-deletion:moment}]
\label{tra:projection-deletion:moment-induction-and-completion}
The rectangle estimate leaves a subexponential regular-trace error. We remove it by separating the low levels, controlled by Proposition~\ref{tra:projection-deletion:level}, from the higher levels, whose regular multiplicities are large. For large \(n\), use \(n^{0.995}\) as the threshold for level. On all types with \(1\le k\le n^{0.995}\), we already have, independently of \eqref{tra:projection-deletion:eq17}, enough decay for the trace estimate \eqref{tra:projection-deletion:eq1} at a sufficiently large absolute power \(p_{\rm low}\). Indeed each \(d_\lambda\le n^k\) by occurrence on injection tuples, and there are at most \(2^k\) diagrams at that level. Using \eqref{tra:projection-deletion:eq3}, for \(n\) sufficiently large and any \(s\ge p_{\rm low}\), the total contribution of these types to \(\operatorname{Tr}|T_n|^{2s}\) is \(o(1)\); for example \(M_0 k/n\le n^{-0.004}\) there, so we can sum the bounds \(2^k n^{2k} (n^{-0.004})^{\alpha k p_{\rm low}}\).

For any type not annihilated by \(T_n\), the number of rows of the diagram is at most \(n/2\). Indeed it needs invariants for the flip subgroup of one dimension, a product of \(n/2\) copies of the two-point symmetric group, so the Pieri rule builds its diagram by adding \(n/2\) horizontal strips. For such diagrams with \(k> n^{0.995}\) and sufficiently large \(n\), we have
\begin{equation}
d_\lambda\ \ge\ \exp(c n^{0.995})
\label{tra:projection-deletion:eq21}
\end{equation}
with absolute \(c>0\). One quick count of tableaux is by ordering the boxes according to increasing sum of their two coordinates, with arbitrary order within each antidiagonal. If there are \(l\) antidiagonals the factorials of their sizes give a product at least \(2^{n-l}\). Here \(l=\max_i(i+\lambda_i-1)\le\sqrt n+\max(\lambda_1,n/2)\) since \(i\lambda_i\le n\), proving \eqref{tra:projection-deletion:eq21}.

When \eqref{tra:projection-deletion:eq17} holds and \(n\) is sufficiently large, \eqref{tra:projection-deletion:eq18} gives
\[
\operatorname{Tr}|T_n|^{2p'}\le \exp(n^{0.99}).
\]
We apply Lemma~\ref{lem:balanced-moment-closure} with
$\nu=2$, $\varepsilon=1/8$ and $P=\max(p_{\rm low},2)$.
The low class consists of the first-row levels
$1\le k\le n^{0.995}$ together with the annihilated blocks;
its contribution is at most $1/16$ above an absolute threshold.
The remaining blocks satisfy \eqref{tra:projection-deletion:eq21}.
Thus the parameters in the common lemma are
\[
 E_n=n^{0.99},\qquad H_n=c n^{0.995},\qquad
 \alpha_d=1+C_3/\log a,\qquad \delta_d=1/\log n.
\]
The child bound $1+1/8$ implies the bound $2$ required in
\eqref{tra:projection-deletion:eq17}, and all size thresholds are
independent of the inherited exponent. Moreover,
\[
 (1+\delta_d)E_n-\delta_dH_n
 =(1+1/\log n)n^{0.99}-c n^{0.995}/\log n\longrightarrow-\infty,
\]
while $\alpha_d(1+\delta_d)=1+O(1/d)$. The lemma supplies bounded
parameters $p_n$ and proves \eqref{tra:projection-deletion:eq1}.

Proposition~\ref{e:nonnormal-moment-conversion}, with singular
moment exponent $2\sup_n p_n$ and remainder $1/8$, now gives the
mixing bound after an absolute number of sweeps. The lower bound in
physical shuffle units is Lemma~\ref{lem:support}.

\end{proof}
